\textbf{Introduction}

According to our current understanding, just after the Big Bang, when the
Universe was extremely hot, the electromagnetic force, the strong nuclear
force as well as the weak nuclear force were unified as one Grand Unified
(GUT) force.

When the Universe cooled down to $T_{GUT} = 10^{29}$\,K, the strong nuclear
force decoupled from the electroweak force. Later, when the temperature
reduced to $T_{EW} = 10^{15}$\,K, the weak force decoupled from the
electromagnetic force. These transitions happened in a rapid succession within
a small fraction of a second after the Big Bang. It is thought that these
phase transitions produced a variety of peculiar objects, called vacuum
defects, which may still be observed today.

This question will discuss properties of one such possible type of defect
called cosmic strings and their observational effects.

\emph{Note 1.} Unless otherwise stated use the laws of Newtonian Mechanics.

\emph{Note 2.} You will use the following constants:
\begin{itemize}
\item Stefan--Boltzmann constant
  \[ \sigma = \frac{2\pi^5 k_B^4}{15\,h^3 c^2}
            = \frac{\pi^2 k_B^4}{60\,\hbar^3 c^2} \]
\item The reduced Planck constant
  \[ \hbar = \frac{h}{2\pi} \]
\item Universal radiation constant
  \[ a = \frac{4\sigma}{c}
       = 7.5657 \times 10^{-16}\,\mathrm{J\,m^{-3}\,K^{-4}} \]
\item Planck temperature
  \[ T_{pl} = \sqrt{\frac{\hbar c^5}{G k_B^2}}
            = 1.416784 \times 10^{32}\,\mathrm{K} \]
\end{itemize}

\emph{Note 3.} Recall that the gravitational field $\vec{g}$ satisfies the
Gauss theorem:
\[ \oint \vec{g}\cdot d\vec{A} = -4\pi G M_{in}, \]
where $M_{in}$ is the mass enclosed by the surface $A$.

\medskip
\textbf{Part A: Gravitational Field of a Cosmic String (22 points).}

As a first approximation, let us consider a cosmic string as an infinitely
long cylinder of radius $r_0$ and mass per unit length $\mu$.

\begin{description}
\item[A.1] Write an expression in terms of the constants $G$, $\mu$ and $r_0$
  for the gravitational field produced by the string, $\vec{g}(r)$. Consider
  the cases $r_0 < r$ and $r_0 > r$ independently. \hfill[6.0\,pt]
\item[A.2] Write an expression in terms of the constants $G$, $\mu$ and $r_0$
  for $g_0 \equiv |\vec{g}(r_0)|$. \hfill[1.0\,pt]
\item[A.3] Let $g$ be defined as $\vec{g}(r)\cdot\hat{r}$. Draw a rough sketch
  of $g$ vs.\ $r$ in the figure given in the answer sheet. \hfill[3.0\,pt]
\item[A.4] It is possible to define a stable orbit around a cosmic string. For
  circular orbits of radius $R > r_0$ and period $\tau$, the following
  relation is attained
  \[ R = A\tau^\alpha, \]
  where $A$ and $\alpha$ are constants. Find $A$ and $\alpha$ in terms of $G$
  and $\mu$. \hfill[4.0\,pt]
\end{description}

The following three questions refer to a classical Newtonian particle moving
with speed $v$ when at a distance $r > r_0$ from the string. You will need to
use the result below:
\[ \int_{x_0}^{x}\frac{dx}{x} = \ln\left(\frac{x}{x_0}\right) \]

\begin{description}
\item[A.5] Show that the gravitational potential energy of the particle is
  \[ U = Gm\mu\ln\left(\frac{r}{b}\right), \]
  where $b$ is any fixed distance. \hfill[3.0\,pt]
\item[A.6] What is the maximum distance, $R_{\max}$, from the string, that the
  particle can reach? \hfill[4.0\,pt]
\item[A.7] Is it possible for the particle to escape the gravitational field?
  Write YES/NO in the answer sheet. \hfill[1.0\,pt]
\end{description}

\medskip
\textbf{Part B: Cosmic string as a photon gas (17 points).}

Consider now a cosmic string as a photon gas inside a very long cylinder of
radius $r_0$ with adiabatic walls, and in thermal equilibrium at temperature
$T$.

\begin{description}
\item[B.1] What is the energy density $\rho$ of the string in terms of $T$,
  $\hbar$, $k_B$ and $c$? \hfill[2.0\,pt]
\item[B.2] The radius $r_0$ is related to the temperature $T$ via
  \[ r_0 = \frac{\hbar^{n_1} c^{n_2}}{k_B T}, \]
  where $\hbar$ is the reduced Planck constant, $c$ is the speed of light in
  vacuum, $k_B$ is the Boltzmann constant, and $n_1$ and $n_2$ are integer
  numbers. Determine $n_1$ and $n_2$. \hfill[4.0\,pt]
\item[B.3] What is the mass per unit length, $\mu$, of the string in terms of
  $\rho$ and $r_0$? \hfill[2.0\,pt]
\item[B.4] Express the inequality for the weak field condition, defined as
  \[ \frac{2G\mu}{c^2} \ll 1, \]
  only in terms of $T$ and $T_{pl}$. \hfill[5.0\,pt]
\item[B.5] Calculate $\frac{2G\mu}{c^2}$ for
  \begin{itemize}
  \item $T = T_{EW}$
  \item $T = T_{GUT}$
  \end{itemize}
  \hfill[3.0\,pt]
\item[B.6] Does the weak field condition hold for $T_{EW}$? Answer YES or NOT.
  Does the weak field condition hold for $T_{GUT}$? Answer YES or NOT.
  \hfill[1.0\,pt]
\end{description}

\medskip
\textbf{Part C: Gravitational Lensing from cosmic Strings (16 points).}

So far, in parts A and B, we have neglected the internal pressure of the
photon gas inside the string. If we include it in our analysis, we need to
consider the General Theory of Relativity.

After solving the Einstein field equations, one finds that the spacetime
around a cosmic string is conical, as if a narrow wedge were removed from a
flat sheet and the edges connected, as shown below.

\ioaafig{2021_TH_Q15-f4.pdf}

{\centering\small
http://www.ctc.cam.ac.uk/outreach/origins/cosmic\_structures\_five.php\par}

A remarkable result of this model is light deflection by a cosmic string,
which leads to the possibility of detection through gravitational lensing.

The angle of deflection (in radians) of a light ray coming from a distant
quasar ($O$ in the figure below), as the light passes close to a cosmic string
($S$ in the figure below) and eventually reaching an observer on the Earth
($E$ in the figure below), is
\[ \delta\phi = \frac{4\pi G\mu}{c^2} \]
and is independent of the parameter, $p$, as shown in the figure below.

In the figure $E$ and $O$ are in a plane perpendicular to the string. The
distance between the observer and the string is $D_{ES}$ and the distance
between the observer and the source is $D_{OE}$.

\ioaafig{2021_TH_Q15-f5.pdf}
% source O offset perpendicular by distance p above the line at the far end
% (source paper Q15-5, below the deflection-angle equation)

\begin{description}
\item[C.1] Although the angle of deflection does not depend on parameter $p$,
  an Earth-based observer will be able to see more than one image only if the
  value of $p$ is within a certain range. Find a condition on the value of the
  parameter $p$ in terms of $D_{ES}$, $D_{OE}$, and temperature $T$, for an
  Earth-based observer to see more than one image of the object $O$.
  \hfill[6.0\,pt]
\item[C.2] In case the observer sees more than one image, what is the angular
  separation between each pair? Find an expression in terms of $D_{ES}$,
  $D_{OE}$ and $\delta\phi$. \hfill[6.0\,pt]
\item[C.3] If $D_{OE} = 2D_{ES}$, determine the minimum size of an optical
  telescope needed to resolve this lensing event produced by a GUT string.
  \hfill[4.0\,pt]
\end{description}
